\chapter{Conclusiones y Líneas Futuras}
\label{ch:conclusiones}

% ---------------------------------------------------------------------------
% SECCIÓN 7.1 — Conclusiones
% ---------------------------------------------------------------------------
\section{Conclusiones}
\label{sec:conclusiones}

El trabajo se propuso construir un sistema de recomendación musical que atendiera al contenido de las letras, y no solamente a las propiedades sonoras de las canciones, con el propósito de aprovechar una dimensión de la obra que los sistemas de uso extendido dejan en un segundo plano. Para la evaluación efectiva, se comprobó que las dos dimensiones consideradas, la sonora y la lírica, organizan el catálogo de manera distinta y aportan información que no se superpone, lo que justificaba combinarlas. Se comprobó luego que esa combinación recomienda mejor que cualquiera de las dos por separado, y que sus sugerencias alcanzan a casi todo el catálogo en lugar de concentrarse en las obras más conocidas. Se comprobó, por último, examinando recomendaciones concretas una por una, que ambas dimensiones cumplen funciones complementarias. El problema que originó el trabajo, incorporar la letra como criterio efectivo de recomendación, se resolvió mediante un sistema verificable.

% ---------------------------------------------------------------------------
% SECCIÓN 7.2 — Líneas Futuras
% ---------------------------------------------------------------------------
\section{Líneas Futuras}
\label{sec:futuras-lineas}

El sistema desarrollado deja planteados varios desarrollos derivados, ordenados de menor a mayor envergadura.

\begin{itemize}[leftmargin=2em]
    \item \textbf{Validación con oyentes}. Someter una muestra de recomendaciones a evaluación de un grupo de personas.

    \item \textbf{Independencia de la fuente de atributos musicales}. Obtener los descriptores sonoros directamente del audio, lo que libera al sistema de una dependencia externa y habilita la ampliación del catálogo.

    \item \textbf{Evolución del prototipo hacia un servicio con retroalimentación}. Registrar la aceptación o el descarte de cada sugerencia, de modo que el sistema incorpore el comportamiento de los oyentes junto al contenido de las obras.
\end{itemize}
